\scanpage{171}
\mspage{87}
% archive sheet number: 170
% line 1
Résumons ce que nous avons obtenu:

% lines 2-12
\uline{Lemme} 4.4.7. Soit $C$ une \inserted{petite} catégorie, $\pi$ un cardinal infini, \deleted{\unsure{supposons}} $J$ une \deleted{petite} \deleted{\unsure{par}} \deleted{\unsure{type}} catégorie \deleted{de cardinal} telle que $\operatorname{card}\operatorname{Fl}J\leq\pi$, $\varphi:J\longrightarrow\operatorname{Ind}'_{\pi}(C)$ un foncteur. Supposons \inserted{les $\varinjlim$ de type $J$ existent dans $C$,} que les \deleted{\unsure{cat.}} foncteurs $\unsure{\varepsilon_j^{\varphi}|\xi_j^{\varphi}}$ ($j\in\operatorname{Ob}J$) de (4.4.6.3) soient \uline{cofinaux} (ce qui s'explicite de façon toute \unsure{simple} si source et but sont filtrants, chose acquise si $C$ est stable par $\varinjlim$ finies). Alors $\varphi$ a une ``représentation indicielle'' canonique \deleted{\unsure{par un $I$}} par
\[
\psi:I\times J\longrightarrow C\tag*{(4.4.7.1)}
\]
avec
\[
I=\underline{\operatorname{Hom}}(J,C)_{/\varphi},\tag*{(4.4.7.1)}
\]
% lines 15-17
i.e. on a l'isom. fonctoriel en $j\in\operatorname{Ob}J$
\[
\varphi\deleted{\unsure{(j)}}\simeq\varinjlim_I\psi(i,j)\tag*{(4.4.7.3)}
\]
\margin{lim. inductive dans $\operatorname{Ind}'_{\pi}(C)$ d'objets de $C$.}

% lines 18-19
Posons pour tout $i\in\operatorname{Ob}I$
\[
\Psi(i)=\varinjlim_J\psi(i,j)\tag*{(4.4.7.4)}
\]
\margin{limite inductive prise dans $C$, mais valable aussi dans $\operatorname{Ind}'_{\pi}C$}

% final lines, continuation next scan
Alors \deleted{$\varphi:J\longrightarrow\operatorname{Ind}'_{\pi}C$ a une} $\varinjlim_{j\in J}\varphi(j)$ (dans $\operatorname{Ind}'_{\pi}C$) existe sss
